% M1_NnaemekaNnachiEgwu.tex
% Milestone 1: Technical Understanding and Research Framing
% Real-Time Systems Seminar, HSHL, Prof. Dr. Stefan Henkler
% Author: Nnaemeka Nnachi-Egwu
%
% NOTE: Per seminar checklist, M1-M3 use structured notes format (bullets,
% formulas, tables, annotations) inside the IEEEtran template shell.
% Publication-style prose is NOT used here (not expected, not appropriate per checklist).

\documentclass[conference,onecolumn]{IEEEtran}
\usepackage{amsmath, amssymb, booktabs, array, microtype, url}

\begin{document}

\title{Milestone 1: Technical Understanding\\
  {\large Priority Inheritance Protocol (Resource Access Protocols)}}

\author{\IEEEauthorblockN{Nnaemeka Nnachi-Egwu}
\IEEEauthorblockA{Electronic Engineering\\
Hochschule Hamm-Lippstadt\\
Lippstadt, Germany\\
nnaemeka.nnachi-egwu@stud.hshl.de}}

\maketitle

%---------------------------------------------------------------
\section{Problem Statement and Motivation}

\textbf{Problem:} Shared exclusive resources in real-time systems cause
\emph{priority inversion}: a high-priority task $\tau_i$ can be blocked
for an \emph{unbounded} duration by lower-priority tasks that do not
even use the contested resource.

\textbf{Formal setting} (Buttazzo \S7.3, p.~209):
\begin{itemize}
  \item $n$ periodic tasks $\{\tau_1, \ldots, \tau_n\}$, priority $P_i$
        (lower number = higher priority, RM assignment).
  \item $m$ shared resources $\{R_1, \ldots, R_m\}$ protected by binary
        semaphores $\{S_1, \ldots, S_m\}$.
  \item Mutual exclusion via \texttt{wait}($S_k$) / \texttt{signal}($S_k$).
  \item Critical section $z_{i,k}$ of $\tau_i$ on $S_k$, duration
        $\delta_{i,k}$; collection $Z_{i,k}$, $\sigma_i$, $\gamma_i$.
\end{itemize}

\textbf{Why it matters (Kopetz definition, RTS slide 7):}
A real-time system is correct iff results are produced at the right
\emph{physical instant}, not only logically.
Priority inversion can cause $\tau_1$ to miss its deadline even with
$U < 1$ (slide 48: ``EDF is optimal only when no resource competition'').
A missed hard deadline is a \emph{correctness failure}, not a
performance degradation.

\textbf{Scenario (Fig.~7.4, Buttazzo p.~208):} $\tau_3$ holds $S_R$;
$\tau_1$ (high prio) arrives and blocks; $\tau_2$ (medium prio, no $S_R$)
preempts $\tau_3$ --- $\tau_1$ blocked while lower-priority $\tau_2$ runs.
Inversion window: $[t_3, t_6]$. Unbounded if multiple $\tau_2$-class
tasks present.

%---------------------------------------------------------------
\section{Core Technical Idea}

PIP (\textbf{Priority Inheritance Protocol}), Sha, Rajkumar, Lehoczky
[SRL90], Buttazzo \S7.6, pp.~214--225:

When $\tau_i$ blocks on $S_k$ held by $\tau_j$ ($P_j < P_i$),
$\tau_j$ \textbf{inherits} $\tau_i$'s priority:
\begin{equation}
  p_j(R_k) = \max\!\bigl\{P_j,\;\max_h\{P_h \mid \tau_h\text{ blocked on }R_k\}\bigr\}.
  \label{eq:pip}
\end{equation}

\textbf{Effect:} Medium-priority tasks (like $\tau_2$) cannot preempt
the boosted $\tau_j$ --- priority inversion bounded.

\textbf{Restoration:} On \texttt{pi\_signal}($S_k$), $\tau_j$ reverts to
$\max\{P_j, \max\{p_h \mid \tau_h\text{ still blocked by }\tau_j\}\}$.

\textbf{Transitivity (Buttazzo p.~215, Lemma~7.2):} If $\tau_3 \to \tau_2
\to \tau_1$ in blocking chain, $\tau_3$ inherits $P_1$ via $\tau_2$.
\emph{Only occurs with nested critical sections.}

%---------------------------------------------------------------
\section{Relevant Formalism}

\subsection{Notation (\S7.3, p.~209)}
\begin{itemize}
  \item $P_i$: nominal priority; $p_i$: active priority
  \item $z_{i,k}$, $Z_{i,k}$: CS and set of CS of $\tau_i$ on $S_k$
  \item $\delta_{i,k}$: CS duration; $\sigma_i$: \# CS of $\tau_i$;
        $\gamma_i$: set of CS that can block $\tau_i$
  \item $l_i$: \# lower-prio tasks that can block $\tau_i$;
        $s_i$: \# semaphores that can block $\tau_i$
  \item $B_i$: worst-case blocking time of $\tau_i$
\end{itemize}

\subsection{Lemmas and Theorems (all Buttazzo \S7.6)}
\begin{itemize}
  \item \textbf{Lemma 7.1} (push-through, p.~218): $S_k$ causes push-through
        blocking to $\tau_i$ iff $\exists\tau_h: P_h > P_i$ using $S_k$
        AND $\exists\tau_j: P_j < P_i$ using $S_k$.
  \item \textbf{Lemma 7.2} (transitivity, p.~218): Transitivity occurs
        only with nested CS.
  \item \textbf{Lemma 7.3} (p.~219): $\tau_i$ blocked at most $l_i$ times.
  \item \textbf{Lemma 7.4} (p.~219): $\tau_i$ blocked at most $s_i$ times.
  \item \textbf{Theorem 7.2} (p.~219, [SRL90]): $\tau_i$ blocked at most
        $\alpha_i = \min(l_i, s_i)$ CS.
  \item \textbf{Lemma 7.5} (p.~221): $Z_{j,k}$ blocks $\tau_i$ only if
        $P_j < P_i \leq C(S_k)$.
\end{itemize}

\subsection{Blocking Time Computation (\S7.6.3, Eqs.~7.9--7.11)}

\[
  B_i^l = \text{sum along task dimension};\quad
  B_i^s = \text{sum along semaphore dimension};\quad
  B_i = \min(B_i^l, B_i^s).
\]

\[
  \gamma_i = \{Z_{j,k} \mid P_j < P_i \leq C(S_k)\}.
\]

Algorithm (Fig.~7.11): complexity $\mathcal{O}(mn^2)$.
\textbf{Not tight} (Buttazzo p.~222): $B_i^l$ can sum CS on same
semaphore (impossible per Lemma~7.4).

\subsection{Schedulability (\S7.9)}

\textbf{Liu-Layland (sufficient, Eq.~7.19):}
\[
  \forall i:\;\sum_{h<i}\frac{C_h}{T_h} + \frac{C_i+B_i}{T_i} \leq i\!\left(2^{1/i}-1\right).
\]

\textbf{RTA (exact, Eq.~7.22):}
\[
  R_i = C_i + B_i + \sum_{h<i}\!\left\lceil\frac{R_i}{T_h}\right\rceil C_h
  \quad\text{(iterate to fixed point)}.
\]

%---------------------------------------------------------------
\section{Assumptions and Scope Limitations}

\begin{itemize}
  \item \textbf{Single processor} (explicit in \S7.3, p.~209).
  \item \textbf{Fixed priorities} (RM/DM); PIP not defined for EDF
        (footnote 2, p.~209; Spuri extension exists but not in Ch.~7).
  \item \textbf{WCET known exactly}: $\delta_{i,k}$ values assumed given.
        In practice, WCET estimation is itself a hard problem [WCET].
  \item \textbf{Binary semaphores}: single-unit resources only.
  \item \textbf{No deadlock prevention}: see \S7.6.5, Fig.~7.13.
  \item \textbf{No chained blocking prevention}: $\alpha_i$ can be $> 1$,
        unlike PCP (Theorem~7.4) which bounds to exactly~1.
  \item \textbf{Transparency assumption}: kernel must implement
        \texttt{pi\_wait}/\texttt{pi\_signal}; application code unchanged.
\end{itemize}

%---------------------------------------------------------------
\section{Runtime, Memory, and Algorithm Aspects}

\begin{itemize}
  \item Blocking computation: $\mathcal{O}(mn^2)$ offline
        (Fig.~7.11 algorithm).
  \item Exact algorithm (Rajkumar): $\mathcal{O}(2^n)$ --- impractical
        for large $n$.
  \item \texttt{pi\_wait} at runtime: unbounded chain traversal (bounded
        in practice, e.g.\ Linux \texttt{RTMUTEX\_DEPTH\_LIMIT}=48).
  \item Memory: extra fields in TCB (active priority, blocked-on semaphore,
        semaphore queue) --- $\mathcal{O}(n + m)$ additional state.
  \item Table~7.5 rates PIP ``hard'' to implement --- the only protocol
        so rated.
\end{itemize}

%---------------------------------------------------------------
\section{Five-Protocol Comparison}

{\small
\begin{tabular}{@{}lccccc@{}}
\toprule
Property & NPP & HLP & PIP & PCP & SRP \\
\midrule
Deadlock-free & $\checkmark$ & $\checkmark$ & \textbf{--} & $\checkmark$ & $\checkmark$ \\
Chained block & -- & -- & \textbf{Yes} & -- & -- \\
Transparent & $\checkmark$ & -- & $\checkmark$ & -- & -- \\
Max blockings & 1 & 1 & $\alpha_i$ & 1 & 1 \\
Impl. effort & Easy & Easy & \textbf{Hard} & Medium & Easy \\
Dynamic prio & -- & -- & -- & -- & $\checkmark$ \\
\bottomrule
\end{tabular}}\\
Source: Buttazzo Table~7.5, p.~248.

%---------------------------------------------------------------
\section{Verified Application Example}

\textbf{Task set:} $\tau_1(C{=}1,T{=}6,D{=}6,\text{prio}=1,\text{CS}=1)$,
$\tau_2(C{=}2,T{=}8,D{=}8,\text{prio}=2,\text{no~CS})$,
$\tau_3(C{=}4,T{=}24,D{=}24,\text{prio}=3,\text{CS}=4)$.
Semaphore $S_R$; ceiling $C(S_R)=P_1$.

\textbf{Blocking:} $B_1=4$ (direct), $B_2=4$ (push-through, Lemma~7.1),
$B_3=0$.

\textbf{RTA:} $R_1=5\leq 6$\,\checkmark, $R_2=8\leq 8$\,\checkmark,
$R_3=8\leq 24$\,\checkmark. All schedulable.

\textbf{Trace highlight:} $\tau_3$ inherits $P_1$ at $t=2$; $\tau_2$
push-through blocked; $\tau_3$ CS ends $t=4$; $\tau_1$ blocking~$=2$
(within $B_1=4$ bound).

\textbf{Open questions noted during study:}
(1)~$\delta{-}1$ vs.\ continuous time: Eq.~(7.9) subtracts~1 (discrete artefact);
hand trace uses full $\delta$ --- both valid in their respective contexts.
(2)~Eq.~(7.11) not tight (Buttazzo p.~222); Rajkumar exact algorithm
($\mathcal{O}(2^n)$) beyond scope here.
(3)~Linux \texttt{PTHREAD\_PRIO\_INHERIT} bounds chain depth at 48 hops
--- formal guarantees may not fully hold on the implementation.
(4)~PIP uniprocessor only --- exact multicore failure mode addressed in M2.

\end{document}
